\begin{proposition}[The equianharmonic crossing]\hypertarget{tag-032N}{}\label{push:src:D.4.5}
\StatementRecord{push:src:D.4.5}{proved}{032N}{\SeriesRef{push:section:D.4}, \SeriesRef{fam:src:C.1.1}, \SeriesRef{fam:src:C.2.3}, \SeriesRef{push:src:D.4.4}}

At the equianharmonic point \(\tau=\rho=e^{\pi i/3}\), with the
original period representative \(\mu=(2-\rho)/3\), one has
\[
 \mathcal N(\mu,\rho)=3^{1/4},\qquad
 D_{\rm part}(z_1)=-\frac3{2\pi}\log48.
\]
Consequently the zero locus passes through the coarse point \(t=0\)
exactly at the height
\[
 H_0=\frac3{2\pi}\log48
 =H_*+\frac3{2\pi}\log\frac43.
\]
It has \(D_c(0)<0\) below this height and \(D_c(0)>0\) above it.
In the fixed coarse coordinate the function is \(C^1\) at zero and
\[
 D_{\rm part}(t)=D_{\rm part}(0)
       +\frac7{32\pi}\operatorname{Re}t+O(|t|^{4/3}),\qquad
 \partial_tD_{\rm part}(t)=\frac7{64\pi}+O(|t|^{1/3}).
\]
Thus at \(H=H_0\) the level set is a regular \(C^1\) arc through
zero, tangent to the imaginary axis.
\end{proposition}


\begin{proof}
\begin{enumerate}[label=\arabic*.,leftmargin=*]
\item For a fixed modulus write \(T=\operatorname{Im}\tau\) and
\[
 u(z)=\log\mathcal N(z,\tau)
 =\log\left|\frac{\vartheta_1(z,\tau)}{\eta_D(\tau)}\right|
       -\frac{\pi(\operatorname{Im}z)^2}{T}.
\]
The theta transformation law makes \(u\) periodic for
\(\Lambda=\mathbb Z+\mathbb Z\tau\). It is locally integrable,
with \(u(z)=\log|z|+c_\tau+o(1)\) at zero, and has no other
singularity on the torus.  In the Jacobi product, with
\(Q=e^{2\pi i\tau}\) and \(X=e^{2\pi iz}\), the quotient is
\[
 \frac{\vartheta_1(z,\tau)}{\eta_D(\tau)}
 =iQ^{1/12}X^{-1/2}(1-X)
       \prod_{m\ge1}(1-Q^mX)(1-Q^mX^{-1}),
\]
up to a constant of absolute value one. For
\(z=x+iy\), \(0<y<T\), every product factor has horizontal
mean logarithm zero, by the power series for
\(\log(1-ae^{2\pi ix})\), \(|a|<1\). The series converges
uniformly on closed substrips. Thus
\[
 \int_0^1u(x+iy)\,dx
       =-\frac{\pi T}{6}+\pi y-\frac{\pi y^2}{T}.
\]
The logarithmic singularities at the strip endpoints are
integrable; alternatively one integrates first on closed substrips
and passes to the limit. Integration in \(y\) gives
\(\int_{\mathbb C/\Lambda}u\,dA/T=0\). A tilted fundamental
parallelogram has the same horizontal averages, by period one.
This proves the zero-mean normalization directly.

\item For \(\tau=\rho\) put \(\alpha=1+\rho\). Since
\(\rho^2=\rho-1\), multiplication by \(\alpha\) on
\((1,\rho)\) has integral matrix
\[
 \begin{pmatrix}1&-1\\1&2\end{pmatrix},\qquad \det=3.
\]
Moreover \(\alpha\mu=1\), \(3\mu=2-\rho\in\Lambda\),
and \(\mu\notin\Lambda\). Its kernel on the torus is therefore
\(K=\{0,\mu,-\mu\}\), and \(|\alpha|^2=3\).
The distributional Euclidean Laplacian is
\[
 \Delta u=2\pi\delta_0-\frac{2\pi}{T}.
\]
Indeed the holomorphic theta factor contributes its simple zero,
and the Gaussian contributes the displayed constant. Pulling back
by the covering \(z\mapsto\alpha z\) shows that
\(u(\alpha z)\) and \(\sum_{v\in K}u(z+v)\) have the same
Laplacian: both have one logarithmic mass at every point of \(K\)
and constant term \(-6\pi/T\). Their difference is consequently
a periodic harmonic function. The logarithmic singularities cancel
locally, so the removable-singularity theorem gives an ordinary
harmonic function on the compact torus. It is constant.
Multiplication by \(\alpha\) and translation preserve normalized
Haar measure. The mean just computed therefore makes that constant
zero, proving the exact distribution identity
\[
 \mathcal N(\alpha z,\rho)
       =\prod_{v\in K}\mathcal N(z+v,\rho).
\]
As \(z\to0\), divide by \(\mathcal N(z,\rho)\). The left
ratio tends to \(|\alpha|=\sqrt3\). Since \(\mathcal N\) is
even, the other two factors tend to \(\mathcal N(\mu,\rho)^2\).
Positivity now gives \(\mathcal N(\mu,\rho)=3^{1/4}\).

\item Finally the exact \cref{fam:src:C.1.1}, \cref{fam:src:C.2.3} normalization is
\[
 D_{\rm part}=\frac6\pi
          \log\frac{\mathcal N(\mu,\tau)}{2\sqrt3}.
\]
Substitution gives \(-3\log48/(2\pi)\). The boundary height is
\(H_*=3\log6/\pi\), so the stated difference is
\(3\log(4/3)/(2\pi)>0\). Addition of \(\operatorname{Im}c\)
proves the crossing height and sign assertions.

\item For regularity in the coarse coordinate, multiplication by \(\rho^2\) is an
isometry of the torus and preserves the zero-mean Green norm just
used. Also \((\rho^2-1)\mu=\rho-1\in\Lambda\). Hence its
nontrivial rotation fixes \(\mu\), and the real gradient of
\(\log\mathcal N\) vanishes there. In the notation of \ref{push:src:D.4.4} this
says \(Z(\rho)=0\). Write \(P_0=\wp'_\rho(\mu)\ne0\),
\(E=\pi^2E_2/3\), and \(b=\operatorname{Im}\mu/T\).
The identities of \ref{push:src:D.4.4} give \(\mu'=b\) at \(\rho\), since
\(g_2(\rho)=0\) and \(L=-2\pi ib\). The heat equation, with
\(\mu\) held fixed, gives
\[
 L_\tau=\frac{-P-2EL}{4\pi i},\qquad L_\mu=-E.
\]
When differentiating \(Z\) along \(\mu(\tau)\), its two terms
involving \(E b\) cancel and the derivative of \(b\) vanishes at
\(\rho\). Thus
\[
 \partial_\tau Z(\rho)=\frac{iP_0}{4\pi},\qquad
 \partial_{\bar\tau}Z(\rho)=0.
\]
Ramanujan's identity and the Weierstrass equation yield
\[
 g_2'(\rho)=-\frac{3i}{\pi}g_3(\rho)
            =\frac{3i}{\pi}P_0^2.
\]
For \(v=\tau-\rho\), it follows that
\[
 3ZP-2g_2=-\frac{21iP_0^2}{4\pi}v+O(|v|^2),\qquad
 \partial_\tau D_{\rm part}
       =\frac{21iP_0^2}{64\pi^4}v^2+O(|v|^3).
\]
Here all remainders are real analytic in \(v,\bar v\).
The exact normalization \(t=j/1728\) gives
\[
 t=-\frac{i g_3(\rho)}{\pi^3}v^3+O(v^4)
   =\frac{iP_0^2}{\pi^3}v^3+O(v^4).
\]
For example this follows immediately from
\(t=E_4^3/(E_4^3-E_6^2)\),
\(E_4'(\rho)=-(2\pi i/3)E_6(\rho)\), and
\(g_3=8\pi^6E_6/27\). Dividing the derivative expansions on
\(v\ne0\) proves
\(\partial_tD_{\rm part}=7/(64\pi)+O(|t|^{1/3})\).
The conjugate derivative has the conjugate expansion. The resulting
real Taylor expansion extends the coarse function as a \(C^1\)
function with nonzero real gradient, giving the asserted local arc
by the implicit function theorem. In particular this is not a
singular level at \(t=0\), despite the vanishing modulus derivative
at the order-three ramification point.

\end{enumerate}
\end{proof}
